\clearpage\begingroup\NoSourceHead\SourceLeaf{484}
\makeatletter\tagsleft@true\let\veqno\@@leqno\makeatother
\renewcommand{\thefootnote}{\arabic{footnote})}\setcounter{footnote}{39}
\ArithmeticHead{27}{R}{GAUSS’ PRINCIPIA GENERALIA THEORIAE FIGURAE FLUIDORUM.}
\noindent\hspace*{\parindent}13. Indem wir die Besprechung der weiteren Umformung von $\delta U$ auf
später verschieben, gehen wir zur Berechnung der \textls[100]{Variationen der Integrale}
\begin{equation*}
s,\ \int z\,ds,\ T\tag{32}
\end{equation*}
über.

Hier ist die von \textsc{Gauss} benutzte Methode (Art. 27) in noch höherem
Grade als bei der Berechnung von $\delta U$ auf das spezielle vorliegende Problem
zugeschnitten. Der Zuwachs, den das Volumen $s$ beim Übergang von der
Fläche $U$ zur variierten Fläche $\bar U$ erfährt, ist das Volumen des zwischen den
beiden Flächen $U$, $\bar U$ einerseits und der Gefässwand andererseits eingeschlossenen
Raumes, vorausgesetzt, dass die verschiedenen Teile desselben mit geeignetem
Vorzeichen in die Rechnung eingestellt werden. Dieser Raum wird nun
folgendermassen in Elemente zerlegt: Entsprechende Punkte der Seiten der
beiden in Nr. 11 eingeführten Dreiecke $MN_1N_2$ und $\bar M\bar N_1\bar N_2$ werden durch
Gerade verbunden. Der so entstehende Körper kann dann angenähert durch
ein Prisma ersetzt werden, dessen mit richtigem Vorzeichen genommenes Volumen
durch das Produkt aus der Basis $dU$ und der Projektion des Variationsvektors
$M\bar M$ auf die äussere Normale der Fläche $U$ im Punkt $M$ ausgedrückt
wird. Dieses Prisma wird als Element des fraglichen Volumens gewählt,
woraus sich für $\delta s$ der Wert
\begin{equation*}
\delta s=\int dU\mathbin{.}\delta e\mathbin{.}\cos(4,5)\tag{33}
\end{equation*}
ergibt, wo
\[
\delta e=\sqrt{(\delta x)^2+(\delta y)^2+(\delta z)^2}=|M\bar M|,
\]
und in der \textsc{Gauss}schen Bezeichnung 4 die Richtung der äusseren Normale an
die Fläche $U$ im Punkt $M$, 5 die Richtung $M\bar M$ bedeutet.

Dieselbe Zerlegung liefert
\begin{equation*}
\delta\int z\,ds=\int z\,dU\mathbin{.}\delta e\mathbin{.}\cos(4,5).\tag{34}
\end{equation*}

In ähnlicher Weise wird die Variation der Fläche $T$ direkt berechnet.
Bezeichnet $P$ die gemeinsame Begrenzungskurve der beiden Flächen $T$ und
$U$, so muss nach der Natur der Aufgabe die Begrenzungskurve $\bar P$ der Fläche
$\bar U$ ebenfalls auf der Gefässwand liegen. Ein Element $dP$ der Kurve $P$, das
ihm mittels (29) zugeordnete Element $d\bar P$ der Kurve $\bar P$ und die beiden, entsprechende
\par\smallskip\noindent\hfill{\small 4\textsuperscript{*}}\par
\clearpage\SourceLeaf{485}
\ArithmeticHead{28}{L}{BOLZA, GAUSS UND DIE VARIATIONSRECHNUNG.}
\noindent Endpunkte der beiden Elemente verbindenden Geraden bilden
dann ein annähernd als Parallelogramm zu betrachtendes Element des Zuwachses,
den die Fläche $T$ bei der Variation erfährt. Bezeichnet daher 8 diejenige
vom Ausgangspunkt des Elementes $dP$ ausgehende Richtung, die auf
diesem Element senkrecht steht, in der Tangentialebene an die Gefässwand
gelegen und in Beziehung auf den Raum $s$ nach aussen gerichtet ist, so ergibt
sich für die Variation von $T$
\begin{equation*}
\delta T=\int dP\mathbin{.}\delta e\mathbin{.}\cos(5,8).\tag{35}
\end{equation*}

Diese \textsc{Gauss}sche Methode der Berechnung der Variationen der drei Integrale
(32) ist an Eleganz und Einfachheit wohl kaum zu übertreffen. Allein
die Schwierigkeiten der Integralrechnung einerseits und der Variationsrechnung
andererseits sind hier so unauflöslich verflochten, dass es schwer sein dürfte,
einen den heutigen Anforderungen an Strenge entsprechenden Beweis darauf
zu gründen, und die Methode ist so ausschliesslich dem speziellen Problem
angepasst, dass es nicht leicht scheint, daraus Folgerungen allgemeiner Natur
für die Variationsrechnung abzuleiten.

Abschliessend können wir die \textsc{Gauss}sche Methode der Berechnung der
ersten Variation eines Doppelintegrals als eine energische Reaktion gegen den
alles überwuchernden Formalismus des $\delta$-Algorithmus, als eine Rückkehr vom
»formalen« zum »gegenständlichen« Denken charakterisieren. Wohl hatte die
Einführung des Variationsprozesses durch \textsc{Lagrange} gegenüber den mühsamen
Infinitesimalbetrachtungen von \textsc{Euler} und seinen Vorgängern einen ganz ungeheuern
Fortschritt bedeutet; aber in der rein mechanischen Anwendung dieses
Prozesses bei gleichzeitiger Variation der unabhängigen Variabeln lag eine
Überschreitung der Grenzen seiner Gültigkeit, die zu Unklarheiten und Fehlern
geführt hatte. \textsc{Poisson} hatte zwar durch die Methode der Parameterdarstellung
die Schwierigkeiten beseitigt, aber nur indem er die Variation der unabhängigen
Variabeln überhaupt über Bord warf. Im Gegensatz dazu hat \textsc{Gauss}
die Methode der Variation der unabhängigen Variabeln wieder zu Ehren gebracht,
indem er die »formale Variation« durch eine sinngemässe, auf die Bedeutung
der dem Formalismus zugrunde liegenden Operation aufgebaute Methode
ersetzte\begingroup\renewcommand{\thefootnote}{39a)}\footnote{Vgl. hierzu die folgende Stelle aus einem Brief von \textsc{Gauss} an \textsc{Schumacher} vom 29. Januar 1829: »\textsc{Lagrange}, wie fast alle Analysten der neueren Zeit, trifft zuweilen der Vorwurf, beim Spiel der Zeichen nicht \textls[100]{immer} die Sache lebendig gegenwärtig zu haben«, \textit{Briefwechsel zwischen C. F. Gauss und H. C. Schumacher}, Bd. II, S. 200.}\endgroup\setcounter{footnote}{39}.
\par\clearpage\SourceLeaf{486}
\ArithmeticHead{29}{R}{GAUSS’ PRINCIPIA GENERALIA THEORIAE FIGURAE FLUIDORUM.}
\begin{center}\textls[100]{B) Umformung der ersten Variation durch partielle Integration.}\end{center}
\noindent\hspace*{\parindent}14. Wir kehren jetzt zum Flächenintegral (30) zurück, durch das die
Variation $\delta U$ dargestellt worden war. Zum Zweck der Umformung desselben
mittels partieller Integration zerlegt \textsc{Gauss} (Art. 22, 23) die Fläche $U$, im
Gegensatz zu der in Nr. 11 benutzten Zerlegungsart, auf eine zweite Art in
Elemente, nämlich durch zwei Systeme von Ebenen, das erste senkrecht zur
$y$-Axe, das zweite senkrecht zur $x$-Axe. Dieser Zerlegung entsprechend wird
der Ausdruck für das Flächenelement
\begin{equation*}
dU=\frac{dx\,dy}{\zeta},\tag{36}
\end{equation*}
wenn $dx,dy$ als positiv angenommen werden und überdies vorausgesetzt wird,
dass auch $\zeta$ auf der ganzen Fläche positiv ist.

Das nach Einführung dieses Ausdrucks für $dU$ in (30) sich ergebende
Doppelintegral für $\delta U$ kann nunmehr mittels der der partiellen Integration
zugrunde liegenden Formeln (2) auf die Form gebracht werden\footnote{Die Grössen $A,B$ entsprechen hier den Grössen $A+V\delta x$, $B+V\delta y$ von Gleichung (22).}
\begin{equation*}
\delta U=\iint\left(\frac{\partial A}{\partial x}+\frac{\partial B}{\partial y}\right)dx\,dy+\iint C\,dx\,dy\tag{37}
\end{equation*}
worin $A,B,C$ gewisse homogene lineare Funktionen von $\delta x$, $\delta y$, $\delta z$ sind.

Hier nimmt nun vor allem die weitere Umformung des ersten Integrals
\begin{equation*}
\iint\left(\frac{\partial A}{\partial x}+\frac{\partial B}{\partial y}\right)dx\,dy\tag{38}
\end{equation*}
unser Interesse in Anspruch. Wir haben früher (Nr. 1) gesehen, wie schon
\textsc{Euler} (1770) an dieser Stelle stecken geblieben war, da aus der Form
\[
\int A\,dy+\int B\,dx,
\]
in die er das Doppelintegral brachte, für die Randbedingung nichts zu
schliessen war, und wie auch \textsc{Lagrange} (1811) und \textsc{Poisson} (1816) hierüber
nicht hinausgekommen waren.

Wir haben dann weiter gesehen (Nr. 4), wie \textsc{Ohm} (1825) eine erste Lösung
\par\clearpage\SourceLeaf{487}
\ArithmeticHead{30}{L}{BOLZA, GAUSS UND DIE VARIATIONSRECHNUNG.}
\noindent der hier vorliegenden Schwierigkeit in der Formel (11) gefunden hatte, sahen
aber gleichzeitig, dass diese Lösung gerade für die bei geometrischen Aufgaben
fast einzig in Betracht kommende Gestalt des Integrationsbereichs illusorisch
wird.

Nunmehr finden wir endlich bei \textsc{Gauss} die endgültige Lösung, die darin
besteht, dass er das Doppelintegral (38) in ein \textls[100]{Linienintegral} verwandelt,
zwar nicht in ein Linienintegral über die Begrenzung des Integrationsbereichs
$\mathfrak{A}$, wie wir es heutzutage mittels der \textsc{Green}schen Formel zu tun gewohnt
sind, vielmehr in ein Linienintegral, genommen über die Begrenzung $P$ des
Flächenstücks $U$. Im einzelnen verfährt dabei \textsc{Gauss} folgendermassen: In dem
Integral
\[
\iint\frac{\partial A}{\partial x}\,dx\,dy
\]
wird zuerst nach $x$ integriert, also die Summation über diejenigen Elemente
der Fläche vorgenommen, welche zwischen zwei Ebenen $y$ und $y+dy$ des
ersten der oben erwähnten Ebenensysteme enthalten sind. Die $x$-Koordinaten
der in gerader Anzahl vorhandenen Schnittpunkte der Kurve $P$ mit der Ebene
$y$ seien der Reihe nach $x^0,x',x'',\ldots$ Dann ist das Resultat dieser Summation
\[
[A(x',y)-A(x^0,y)+A(x''',y)-A(x'',y)\ldots]dy.
\]

Sind nun $dP^0$, $dP'$, $dP''$, $\ldots$ die Elemente der Kurve $P$, die zwischen den
beiden Ebenen $y$ und $y+dy$ gelegen sind, und sind $X,Y,Z$ die Richtungskosinus
der positiven Tangente an die Kurve $P$ in einem allgemeinen Punkte
derselben, so ist bei der von \textsc{Gauss} gewählten Festsetzung über die positive
Umlaufsrichtung der Kurve $P$
\begin{equation*}
dy=-Y^0dP^0=+Y'dP'=-Y''dP''\ldots\tag{39}
\end{equation*}

Die obige Summe lässt sich also schreiben
\[
A(x_0,y)Y^0dP^0+A(x',y)Y'dP'+A(x'',y)Y''dP''+\ldots,
\]
und wenn nunmehr auch nach $y$ summiert wird, so erhält man das Resultat
\[
\iint\frac{\partial A}{\partial x}\,dx\,dy=\int AY\,dP,
\]
wo das Linienintegral rechts über die ganze Begrenzungskurve $P$ in positivem
\par\clearpage\SourceLeaf{488}
\ArithmeticHead{31}{R}{GAUSS’ PRINCIPIA GENERALIA THEORIAE FIGURAE FLUIDORUM.}
\noindent Sinn zu erstrecken ist. Das analoge Verfahren, bei dem nunmehr zuerst nach
$y$ integriert wird, ergibt
\[
\iint\frac{\partial B}{\partial y}\,dx\,dy=-\int BX\,dP,
\]
da hier in entsprechender Bezeichnung
\begin{equation*}
dx=+X^0dP^0=-X'dP'=+X''dP''\ldots,\tag{40}
\end{equation*}
und die Verbindung beider Resultate liefert den \textls[100]{Satz}
\begin{equation*}
\iint\left(\frac{\partial A}{\partial x}+\frac{\partial B}{\partial y}\right)dx\,dy=\int(AY-BX)\,dP,\tag{41}
\end{equation*}
wenn $P$ irgend eine geschlossene Raumkurve ist, deren Projektion auf die
$x,y$-Ebene die einmal in positivem Sinn durchlaufene Begrenzungskurve des
Integrationsbereichs $\mathfrak{A}$ des Doppelintegrals ist.

\textsc{Gauss} spricht das zwar nicht als allgemeinen Satz aus mit beliebigen
Funktionen $A,B$, sondern nur für die speziellen Funktionen $A,B$, die bei
der Berechnung von $\delta U$ auftreten. Trotzdem können wir den Satz mit Fug
und Recht als die \textsc{Gauss}sche Form der \textsc{Green}schen Formel bezeichnen,
denn die \textsc{Green}sche Formel für die Ebene ist in (41) als spezieller Fall enthalten,
wenn nämlich die Kurve $P$ eben ist und in der $x,y$-Ebene liegt.
\par\clearpage\endgroup
